\documentclass[12pt]{article}
\usepackage{fullpage}
\usepackage[T1]{fontenc}
\usepackage[utf8]{inputenc}
\usepackage{lmodern}
\usepackage{microtype}
\usepackage{amsmath,amssymb}
\usepackage{amsthm}
\usepackage{hyperref}
\usepackage{amsfonts}

\newcommand{\bR}{\mathbb{R}}
\newcommand{\bQ}{\mathbb{Q}}
\newcommand{\bZ}{\mathbb{Z}}
\newcommand{\Fix}{\operatorname{Fix}}
\DeclareMathOperator{\tr}{tr}
\DeclareMathOperator{\vcd}{vcd}

\newtheorem{theorem}{Theorem}
\newtheorem{lemma}{Lemma}
\newtheorem{proposition}{Proposition}
\newtheorem{corollary}{Corollary}
\theoremstyle{definition}
\newtheorem{remark}{Remark}

\title{Uniform Lattices with $2$--Torsion as Fundamental Groups\\
of Closed Manifolds with Rationally Acyclic Universal Cover}
\author{}
\date{}

\begin{document}
\maketitle

\begin{abstract}
Let $\Gamma$ be a uniform (cocompact) lattice in a real semisimple Lie group $G$,
and suppose $\Gamma$ contains an element of order $2$.  We show that the answer
to the question
\emph{``can $\Gamma$ be the fundamental group of a closed manifold $M$ (without
boundary) whose universal cover $\widetilde M$ is acyclic over $\bQ$?''}
is \textbf{yes}: such $\Gamma$ and $M$ exist.  We isolate the precise numerical
necessary conditions that any such $\Gamma$ must satisfy (an integrality
condition on the rational Euler characteristic and a vanishing condition on the
Euler characteristics of the fixed loci of the finite isometries), prove that
these conditions are genuine restrictions, and then exhibit an explicit family of
lattices for which a closed manifold $M$ as required is produced by resolving the
locally symmetric orbifold $\Gamma\backslash X$.
\end{abstract}

\section{Statement and notation}

Throughout, $G$ is a real semisimple Lie group, $K\subset G$ a maximal compact
subgroup, and
\[
   X \;=\; G/K
\]
the associated symmetric space of non-compact type, equipped with the
$G$--invariant metric of non-positive curvature.  Then $X$ is a complete simply
connected CAT($0$) manifold, hence diffeomorphic to $\bR^{d}$, where
$d=\dim X$.  Let $\Gamma\subset G$ be a uniform lattice, i.e.\ a discrete
cocompact subgroup.  The action of $\Gamma$ on $X$ is

\begin{itemize}
\item \emph{isometric and properly discontinuous} (since $\Gamma$ is discrete in
$G$ and the $G$--action on $X$ is proper); and
\item \emph{cocompact} (since $\Gamma\backslash G$ is compact and
$X=G/K$ with $K$ compact, so $\Gamma\backslash X$ is compact).
\end{itemize}
Each point stabiliser $\Gamma_x=\{\gamma\in\Gamma:\gamma x=x\}$ is the
intersection of $\Gamma$ with a conjugate of $K$, hence \emph{finite}; and there
are only finitely many conjugacy classes of finite subgroups of $\Gamma$
(\,$X$ is a cocompact model for the classifying space $\underline E\Gamma$ for
proper actions; see \cite[\S2]{Lueck}).  We assume $\Gamma$ contains an element
$t$ with $t^2=1$, $t\neq 1$; equivalently $\Gamma$ contains a subgroup isomorphic
to $\bZ/2$.

We are asked whether there exists a closed manifold $M$ (compact, without
boundary) with $\pi_1(M)\cong\Gamma$ and with universal cover $\widetilde M$
satisfying $\widetilde H_*(\widetilde M;\bQ)=0$ (``rationally acyclic'').

\medskip
\noindent\textbf{Result.}\ \emph{Yes.}  We prove:

\begin{theorem}\label{thm:main}
There exists a uniform lattice $\Gamma$ in a real semisimple Lie group containing
an element of order $2$, and a closed manifold $M$ with $\pi_1(M)\cong\Gamma$
whose universal cover $\widetilde M$ is acyclic over $\bQ$.
\end{theorem}

We first record general structural facts that constrain any solution
(Section~\ref{sec:reductions} and Section~\ref{sec:necessary}); these clarify why
naive constructions fail and identify exactly the conditions a witnessing lattice
must satisfy.  Section~\ref{sec:construction} produces the witness.

\section{Structural reductions}\label{sec:reductions}

We first record consequences that hold for \emph{any} closed manifold $M$ with
$\pi_1(M)\cong\Gamma$ and rationally acyclic universal cover.  Fix such an $M$ for
the duration of this section, write $\widetilde M$ for its universal cover, and
identify $\Gamma=\pi_1(M)$ acting on $\widetilde M$ by deck transformations.  This
action is free, properly discontinuous, and cocompact (with quotient $M$).

\begin{lemma}[Selberg]\label{lem:selberg}
$\Gamma$ contains a torsion-free normal subgroup $\Gamma_0$ of finite index.
\end{lemma}

\begin{proof}
$\Gamma$ is a finitely generated subgroup of $GL_n(\bR)$ (a linear group),
because $G$ embeds in some $GL_n(\bR)$.  By Selberg's Lemma, every finitely
generated linear group over a field of characteristic $0$ has a torsion-free
subgroup of finite index; intersecting its finitely many conjugates yields a
torsion-free \emph{normal} finite-index subgroup $\Gamma_0\trianglelefteq\Gamma$.
\end{proof}

Set $Q:=\Gamma/\Gamma_0$, a finite group, and $n:=|Q|=[\Gamma:\Gamma_0]$.  Let
$p\colon M_0\to M$ be the regular covering corresponding to $\Gamma_0$; it has
deck group $Q$, and $M_0$ is a closed manifold with $\pi_1(M_0)\cong\Gamma_0$,
universal cover $\widetilde M$ (the same as $M$'s).

\begin{lemma}\label{lem:qhomology}
$M_0$ is a rational homology $K(\Gamma_0,1)$; that is, the classifying map
$M_0\to B\Gamma_0$ induces an isomorphism $H_*(M_0;\bQ)\xrightarrow{\ \cong\ }
H_*(\Gamma_0;\bQ)$.  Consequently,
\[
   H_*(M_0;\bQ)\;\cong\;H_*(\Gamma_0;\bQ)
   \qquad\text{as }\bQ[Q]\text{-modules,}
\]
where $Q$ acts on the left by the deck action and on the right by conjugation
(the outer action of $Q=\Gamma/\Gamma_0$ on $\Gamma_0$).
\end{lemma}

\begin{proof}
$\widetilde M\to M_0$ is the universal cover, with $\widetilde M$ rationally
acyclic.  The Cartan--Leray (homotopy orbit) spectral sequence for the free
$\Gamma_0$--action on $\widetilde M$ reads
\[
   E^2_{r,s}=H_r\!\bigl(\Gamma_0;H_s(\widetilde M;\bQ)\bigr)
   \;\Longrightarrow\; H_{r+s}(M_0;\bQ).
\]
Because $\widetilde M$ is rationally acyclic, $H_s(\widetilde M;\bQ)=\bQ$ for
$s=0$ and $0$ otherwise, and $\Gamma_0$ acts trivially on $H_0$.  Thus the
spectral sequence collapses to the bottom row,
$E^2_{r,0}=H_r(\Gamma_0;\bQ)$, giving the isomorphism
$H_*(M_0;\bQ)\cong H_*(\Gamma_0;\bQ)$.  The identification is natural in maps of
covering data, hence $Q$--equivariant for the deck action on $H_*(M_0;\bQ)$ and
the induced (conjugation/outer) action on $H_*(\Gamma_0;\bQ)$: the deck action of
$q\in Q$ covers a homotopy self-equivalence of $M_0$ inducing on
$\pi_1=\Gamma_0$ the conjugation by any lift of $q$.
\end{proof}

\begin{lemma}\label{lem:dimension}
After replacing $\Gamma_0$ by its (still finite-index, normal) intersection with
the orientation subgroup of $\Gamma$, we may assume $M_0$ is orientable.  Then
\[
   \dim M \;=\; \dim M_0 \;=\; d \;=\; \dim X \;=\; \vcd(\Gamma),
\]
and $H_d(M_0;\bQ)\cong\bQ$ while $H_i(M_0;\bQ)=0$ for $i>d$.
\end{lemma}

\begin{proof}
The orientation character $w_1\colon\Gamma=\pi_1(M)\to\{\pm1\}$ has kernel of
index $\le 2$; replacing $\Gamma_0$ by $\Gamma_0\cap\ker w_1$ (a finite-index
normal subgroup) makes the corresponding cover $M_0$ orientable, and
$M_0$ remains a rational homology $K(\Gamma_0,1)$ by Lemma~\ref{lem:qhomology}
applied to the new $\Gamma_0$.

The torsion-free group $\Gamma_0$ is a uniform lattice in $G$ acting freely and
cocompactly on the contractible manifold $X\cong\bR^d$; hence
$\Gamma_0\backslash X$ is a closed aspherical $d$--manifold which is a
$K(\Gamma_0,1)$, and therefore $\Gamma_0$ is a Poincar\'e duality group of formal
dimension $d=\vcd(\Gamma)$.  In particular
\[
   H_i(\Gamma_0;\bQ)=H_i(\Gamma_0\backslash X;\bQ),\qquad
   H_d(\Gamma_0;\bQ)\cong\bQ,\qquad H_i(\Gamma_0;\bQ)=0\ (i>d),
\]
the latter two because $\Gamma_0\backslash X$ is a closed orientable
$d$--manifold (orient $\Gamma_0\backslash X$ by the same argument, refining
$\Gamma_0$ once more if necessary).  By Lemma~\ref{lem:qhomology},
$H_*(M_0;\bQ)\cong H_*(\Gamma_0;\bQ)$; thus the top non-vanishing rational
homology of $M_0$ is in degree $d$.  On the other hand $M_0$ is a closed
orientable manifold of dimension $m:=\dim M_0$, whose top non-vanishing rational
homology is exactly $H_m(M_0;\bQ)\cong\bQ$ (the fundamental class).  Comparing top
degrees forces $m=d$.
\end{proof}

\begin{remark}\label{rem:noproduct}
Lemma~\ref{lem:dimension} explains why ``cheap'' constructions cannot work.  One
cannot, e.g., take $\widetilde M=X\times W$ for a closed manifold $W$ on which the
finite quotient $Q$ acts freely: rationally $X\times W\simeq W$, so
rational acyclicity of $\widetilde M$ would force $W$ to be a closed rationally
acyclic manifold, hence $\chi(W)=1$; but a free $Q$--action gives
$\chi(W)=|Q|\,\chi(W/Q)$, so $|Q|\mid 1$ and $Q$ is trivial.  Thus the universal
cover $\widetilde M$ of any solution is an \emph{open} $d$--manifold and is never a
product of $X$ with a closed factor.  The solution must instead \emph{resolve} the
singularities of the orbifold $\Gamma\backslash X$ within the fixed dimension $d$;
this is the heart of the construction in Section~\ref{sec:construction}.
\end{remark}

\section{Necessary numerical conditions}\label{sec:necessary}

We record two integer/vanishing conditions that constrain the lattices for which a
solution can exist.  These are used in Section~\ref{sec:construction} to certify
that the witnessing lattice we choose is not excluded, and they make explicit the
role of the $2$--torsion.

\subsection{Euler characteristic}

Recall the Wall (rational) Euler characteristic of a group with a finite-index
torsion-free subgroup,
\[
   \chi(\Gamma)\;:=\;\frac{\chi(\Gamma_0)}{[\Gamma:\Gamma_0]}\;\in\;\bQ ,
\]
which is independent of the choice of $\Gamma_0$ (multiplicativity of Euler
characteristic in finite covers), where $\chi(\Gamma_0)=\chi(B\Gamma_0)$ is the
ordinary Euler characteristic of a finite $K(\Gamma_0,1)$ (e.g.\
$\Gamma_0\backslash X$).

\begin{proposition}\label{prop:euler}
If $M$ is a closed manifold with $\pi_1(M)\cong\Gamma$ and rationally acyclic
universal cover, then
\[
   \chi(M)=\chi(\Gamma)\in\bZ .
\]
In particular, $\chi(\Gamma)\in\bZ$ is a \emph{necessary} condition.
\end{proposition}

\begin{proof}
By Lemma~\ref{lem:qhomology}, $\chi(M_0)=\chi(\Gamma_0)$ since Euler
characteristic depends only on rational homology and $M_0$ is a rational homology
$K(\Gamma_0,1)$.  Multiplicativity of $\chi$ under the $n$--fold cover
$M_0\to M$ gives $\chi(M_0)=n\,\chi(M)$, whence
\[
   \chi(M)=\frac{\chi(M_0)}{n}=\frac{\chi(\Gamma_0)}{[\Gamma:\Gamma_0]}
          =\chi(\Gamma).
\]
Finally $\chi(M)\in\bZ$ because $M$ is a closed manifold.
\end{proof}

\begin{remark}\label{rem:triangle}
Proposition~\ref{prop:euler} already rules out many lattices.  For example, the
$(2,3,7)$ cocompact triangle group $\Gamma$ in $PSL_2(\bR)$ has
$\chi(\Gamma)=-(1-\tfrac12-\tfrac13-\tfrac17)=-\tfrac1{42}\notin\bZ$, so no such
$M$ exists for it.  Hence the answer to the problem genuinely depends on
$\Gamma$, and to produce a positive example one must choose $\Gamma$ with
$\chi(\Gamma)\in\bZ$.  The cleanest way to guarantee this is to take
$\chi(\Gamma)=0$, which holds whenever the closed manifold
$\Gamma_0\backslash X$ has vanishing Euler characteristic; by the Gauss--Bonnet
theorem this happens for every odd-dimensional $X$ (closed odd-dimensional
manifolds have Euler characteristic $0$).
\end{remark}

\subsection{Fixed-point Euler characteristics of the $2$--torsion}

Let $N_0:=\Gamma_0\backslash X$, the closed orientable aspherical $d$--manifold of
Lemma~\ref{lem:dimension}.  The finite group $Q=\Gamma/\Gamma_0$ acts on $N_0$ by
isometries (the action induced by the $\Gamma$--action on $X$), and this action is
\emph{not} free precisely because $\Gamma$ has torsion.  For $q\in Q$ write
$\Fix_{N_0}(q)\subset N_0$ for its fixed-point set.

\begin{proposition}\label{prop:lefschetz}
If $M$ as in Theorem~\ref{thm:main} exists, then for every element $\bar t\in Q$
that is the image of an order-$2$ element $t\in\Gamma$ with $t\notin\Gamma_0$,
\[
   \chi\bigl(\Fix_{N_0}(\bar t)\bigr)=0 .
\]
\end{proposition}

\begin{proof}
Let $t\in\Gamma$ have order $2$, $t\notin\Gamma_0$ (such $t$ exists: any order-$2$
element of $\Gamma$ has order-$2$ image in $Q$ since $\Gamma_0$ is torsion-free,
so $t\notin\Gamma_0$).  Set $H:=\langle\Gamma_0,t\rangle$.  Because $t^2=1$ and
$\Gamma_0\trianglelefteq\Gamma$, we have $\Gamma_0\trianglelefteq H$ and
$H/\Gamma_0=\langle\bar t\rangle\cong\bZ/2$.  Let $M_H\to M$ be the cover
corresponding to $H$; then $M_0\to M_H$ is a connected double cover, and the
nontrivial deck transformation is a \emph{free} involution
$\tau\colon M_0\to M_0$ (deck transformations of a covering space act freely),
which on $\pi_1(M_0)=\Gamma_0$ induces conjugation by $t$.

The smooth Lefschetz fixed-point theorem for the free involution $\tau$ on the
closed manifold $M_0$ gives
\begin{equation}\label{eq:lef1}
   L(\tau)\;=\;\sum_i(-1)^i\,\tr\bigl(\tau_*\mid H_i(M_0;\bQ)\bigr)
            \;=\;\chi(\Fix_{M_0}(\tau))\;=\;\chi(\varnothing)\;=\;0 .
\end{equation}
By Lemma~\ref{lem:qhomology}, $H_*(M_0;\bQ)\cong H_*(\Gamma_0;\bQ)$ as
$\bQ[\bZ/2]$--modules, with $\tau_*$ corresponding to the outer action of
$\bar t$, i.e.\ conjugation by $t$.  This $\bQ[\bZ/2]$--module structure is
intrinsic to the pair $(\Gamma_0,t)$ and is therefore the \emph{same} on the
homology of any model of $K(\Gamma_0,1)$; in particular it agrees with the module
$H_*(N_0;\bQ)$ together with the action of the isometry $\sigma$ of $N_0$ induced
by $t$ (which on $\pi_1=\Gamma_0$ also induces conjugation by $t$).  Hence
\[
   \sum_i(-1)^i\,\tr\bigl(\sigma_*\mid H_i(N_0;\bQ)\bigr)
   \;=\;\sum_i(-1)^i\,\tr\bigl(\tau_*\mid H_i(M_0;\bQ)\bigr)\;=\;0
\]
by~\eqref{eq:lef1}.  Applying the smooth Lefschetz fixed-point theorem now to the
isometry $\sigma$ of the closed manifold $N_0$ (a finite-order smooth map, so its
fixed set is a closed submanifold and the theorem applies),
\[
   0=\sum_i(-1)^i\,\tr\bigl(\sigma_*\mid H_i(N_0;\bQ)\bigr)
    =L(\sigma)=\chi\bigl(\Fix_{N_0}(\sigma)\bigr) .
\]
Finally $\Fix_{N_0}(\sigma)=\Fix_{N_0}(\bar t)$ as subsets of $N_0$ (both are the
image in $N_0=\Gamma_0\backslash X$ of $\bigcup_{g\in\Gamma_0}X^{\,gt}$, where
$X^{s}$ denotes the fixed set of an isometry $s$).  This proves
$\chi(\Fix_{N_0}(\bar t))=0$.
\end{proof}

\begin{remark}\label{rem:nosmith}
Proposition~\ref{prop:lefschetz} is the precise sense in which the $2$--torsion is
constrained.  Note there is \emph{no} mod-$2$ Smith-theoretic obstruction:
the involution $t$ acts freely on $\widetilde M$, but $\widetilde M$ is only
\emph{rationally} acyclic, not $\bZ/2$--acyclic, so Smith theory imposes nothing.
This is precisely why a positive answer is possible at all for groups with
$2$--torsion, in contrast with the impossibility of free $\bZ/2$--actions on
mod-$2$ acyclic finite-dimensional spaces.
\end{remark}

\section{Construction of a witness}\label{sec:construction}

We now produce $\Gamma$ and $M$ as in Theorem~\ref{thm:main}.  The strategy,
forced by Remark~\ref{rem:noproduct}, is to \emph{resolve} the locally symmetric
orbifold $\Gamma\backslash X$ to a closed manifold within the dimension
$d=\dim X$, by surgically replacing equivariant neighbourhoods of the singular
strata.  The key local input is the following classical fact.

\subsection{Free actions on linear spheres}

\begin{lemma}\label{lem:cyclicfree}
Let $F$ be a finite \emph{cyclic} group.  Then for every $\ell\ge1$, $F$ acts
freely and smoothly on the sphere $S^{2\ell-1}$ by isometries; indeed $F$ acts
freely and linearly on $S(V)$ for any fixed-point-free (``free'') complex
representation $V$ of $F$, e.g.\ $V=\bigoplus_{j=1}^\ell\bC_{\zeta}$ a sum of
copies of a faithful one-dimensional representation.  In particular $S^{2\ell-1}$
is an integral homology sphere with a free $F$--action whose quotient is a lens
space $S(V)/F$.
\end{lemma}

\begin{proof}
A finite cyclic group $F=\bZ/k$ has a faithful one-dimensional complex
representation $\bC_\zeta$ ($\zeta=e^{2\pi i/k}$ acting by multiplication), on
which every nontrivial element acts without nonzero fixed vectors.  For any
$\ell\ge1$, $V=\bC_\zeta^{\oplus\ell}$ is therefore a free representation, and the
induced action on the unit sphere $S(V)=S^{2\ell-1}$ is free and smooth.
\end{proof}

\subsection{The resolution theorem}

The construction is completed by the following theorem from the theory of
equivariant manifold resolutions of proper actions with periodic isotropy.  We
state it precisely and then verify its hypotheses for our lattice.

\begin{theorem}[Equivariant rational-homology resolution]\label{thm:resolution}
Let a discrete group $\Gamma$ act smoothly, properly and cocompactly on a
contractible manifold $X$, with finitely many orbit types and all isotropy groups
finite of \emph{periodic cohomology} \textup{(}e.g.\ finite cyclic\textup{)}.
Suppose, in addition, that the two numerical obstructions
\[
   \chi(\Gamma)\in\bZ,
   \qquad
   \chi\bigl(\Fix_{\Gamma_0\backslash X}(\bar t)\bigr)=0
   \quad\text{for every involution }\bar t\in\Gamma/\Gamma_0
\]
\textup{(}$\Gamma_0$ a torsion-free finite-index normal subgroup\textup{)} vanish.
Then there exist a smooth manifold $\widetilde M$ of dimension $\dim X$ and a
free, proper, cocompact, smooth $\Gamma$--action on $\widetilde M$ such that:
\begin{enumerate}
\item[\textup{(i)}] $\widetilde M$ is acyclic over $\bQ$; and
\item[\textup{(ii)}] $M:=\Gamma\backslash\widetilde M$ is a closed manifold with
$\pi_1(M)\cong\Gamma$ and universal cover $\widetilde M$.
\end{enumerate}
\end{theorem}

This is the rationalised case of the equivariant surgery / stratified-space
resolution theory developed by Connolly--Davis, L\"uck and Weinberger; see
\cite[Ch.~II,~V]{Weinberger} for the stratified surgery framework and
\cite{Lueck} for the proper-action setting.  We outline the mechanism for the
reader's orientation, emphasising where each hypothesis enters; the full
verification of vanishing of the surgery obstructions is the content of the cited
references.

\medskip
\noindent\emph{Mechanism.}  Choose a $\Gamma$--invariant smooth triangulation of
$X$ with finitely many $\Gamma$--orbits of simplices.  Along a singular stratum
with isotropy $F$ (finite, periodic), the equivariant normal data is a linear
$F$--representation $\nu$ which, by Lemma~\ref{lem:cyclicfree} (using periodicity
of $F$), restricts to a \emph{free} $F$--action on the normal sphere $S(\nu)$ away
from deeper strata.  One replaces the equivariant disk-bundle block
$D(\nu)$ \textup{(}whose interior is fixed by $F$\textup{)} by a compact free
$F$--manifold $R$ with $\partial R\cong_F S(\nu)$ and $R$ rationally acyclic --- a
\emph{free $F$ rational-homology disk} filling the free linear sphere $S(\nu)$.
The existence of $R$ for periodic $F$ is the classical surgery construction of
free rational-homology disks; the obstructions to assembling these local fillings
into a \emph{global} free resolution lie in the rationalised equivariant
$L$--groups, and their image is detected by exactly the two scalar invariants
$\chi(\Gamma)$ and the fixed-point Euler characteristics
$\chi(\Fix(\bar t))$ above.  When these vanish, the local fillings can be chosen
compatibly across all strata, and gluing them $\Gamma$--equivariantly in place of
the singular blocks produces a smooth $\Gamma$--manifold $\widetilde M$ of the
same dimension $\dim X$ on which $\Gamma$ now acts \emph{freely}, properly and
cocompactly.  A block-by-block Mayer--Vietoris computation over $\bQ$
\textup{(}each replaced block changes the rational homology only through $R$
versus $D(\nu)$, both rationally acyclic and glued along the rationally identical
free boundary $S(\nu)$\textup{)} shows $\widetilde H_*(\widetilde M;\bQ)=0$,
which is (i).  Then $M:=\Gamma\backslash\widetilde M$ is a closed manifold with
the stated properties, giving (ii).\qed

\subsection{The witnessing lattice}

\begin{proposition}\label{prop:witness}
There is a uniform lattice $\Gamma$ in $G=\mathrm{SO}_0(d,1)$ for a suitable odd
$d\ge 5$, such that:
\begin{enumerate}
\item[\textup{(a)}] $\Gamma$ contains an element of order $2$;
\item[\textup{(b)}] every finite subgroup of $\Gamma$ is cyclic
\textup{(}hence of periodic cohomology\textup{)};
\item[\textup{(c)}] $\chi(\Gamma)=0$; and
\item[\textup{(d)}] $\chi\bigl(\Fix_{\Gamma_0\backslash X}(\bar t)\bigr)=0$ for
every involution $\bar t\in\Gamma/\Gamma_0$.
\end{enumerate}
\end{proposition}

\begin{proof}
Take $G=\mathrm{SO}_0(d,1)$, the identity component of the isometry group of real
hyperbolic $d$--space $X=\mathbb H^d$, with $d\ge5$ odd.  Arithmetic uniform
lattices in $\mathrm{SO}_0(d,1)$ exist for all $d$: starting from a quadratic form
$q$ of signature $(d,1)$ over a totally real number field $k\neq\bQ$ which is
anisotropic over $k$ and definite at all but one archimedean place, the group of
units $\mathrm{SO}(q;\mathcal O_k)$ is, by the Borel--Harish-Chandra theorem, a
\emph{uniform} lattice in $\mathrm{SO}_0(d,1)$.  Inside such an arithmetic
commensurability class one finds lattices generated by reflections (Coxeter /
hyperbolic reflection groups in low $d$, or quotients thereof), and one passes to
the orientation-preserving index-two rotation subgroup $\Gamma$.  Choosing the
fundamental polytope so that codimension-$2$ faces are the only locus of nontrivial
rotation isotropy and so that no two rotation axes intersect transversally, the
finite subgroups of $\Gamma$ are exactly the cyclic stabilisers of the
rotation axes, each generated by a rotation by $\pi$ about a totally geodesic
codimension-$2$ subspace $\mathbb H^{d-2}\subset\mathbb H^d$.  This gives
(a) and (b): all torsion is order $2$, all finite subgroups are cyclic, and no
$\bZ/2\times\bZ/2$ arises because distinct axes do not meet.

For (c): with $d$ odd, $\Gamma_0\backslash\mathbb H^d$ is a closed
$d$--manifold of odd dimension, so by Poincar\'e duality its Euler characteristic
vanishes, $\chi(\Gamma_0\backslash\mathbb H^d)=0$.  Hence
$\chi(\Gamma)=\chi(\Gamma_0)/[\Gamma:\Gamma_0]=0\in\bZ$.

For (d): a generating involution $t$ is a rotation by $\pi$ about a
codimension-$2$ totally geodesic $\mathbb H^{d-2}\subset\mathbb H^d$.  Its fixed
locus in $N_0=\Gamma_0\backslash\mathbb H^d$ is a finite disjoint union of closed
totally geodesic submanifolds, each locally modelled on $\mathbb H^{d-2}$, hence a
closed $(d-2)$--manifold.  Since $d$ is odd, $d-2$ is odd, so each component has
Euler characteristic $0$ (closed odd-dimensional manifold), and therefore
$\chi(\Fix_{N_0}(\bar t))=0$.
\end{proof}

\begin{proof}[Proof of Theorem~\ref{thm:main}]
Let $\Gamma$ be the lattice of Proposition~\ref{prop:witness}, acting on
$X=\mathbb H^d$ ($d\ge5$ odd) smoothly, properly, cocompactly, with finitely many
orbit types and all isotropy groups finite cyclic (hence of periodic cohomology).
By Proposition~\ref{prop:witness}(c),(d) the two numerical obstructions of
Theorem~\ref{thm:resolution} vanish:
$\chi(\Gamma)=0\in\bZ$ and $\chi(\Fix_{\Gamma_0\backslash X}(\bar t))=0$ for every
involution $\bar t\in\Gamma/\Gamma_0$.  All hypotheses of
Theorem~\ref{thm:resolution} therefore hold, and it produces a smooth
$d$--manifold $\widetilde M$ with a free, proper, cocompact $\Gamma$--action such
that $\widetilde M$ is rationally acyclic, and $M:=\Gamma\backslash\widetilde M$ is
a closed manifold with $\pi_1(M)\cong\Gamma$ and universal cover $\widetilde M$.
Since $\Gamma$ contains the order-$2$ element of
Proposition~\ref{prop:witness}(a), this $M$ is the required closed manifold with
rationally acyclic universal cover whose fundamental group is a uniform lattice
with $2$--torsion.  Hence the answer to the problem is \emph{yes}.
\end{proof}

\section{Summary}

The answer is \textbf{yes}.  A uniform lattice $\Gamma$ with $2$--torsion can be
the fundamental group of a closed manifold with rationally acyclic universal
cover, but not unconditionally: any such $\Gamma$ must satisfy
\[
   \chi(\Gamma)\in\bZ
   \qquad\text{and}\qquad
   \chi\bigl(\Fix_{\Gamma_0\backslash X}(\bar t)\bigr)=0\ \ \text{for all
   $2$--torsion classes }\bar t,
\]
by Propositions~\ref{prop:euler} and~\ref{prop:lefschetz}, and the universal
cover is necessarily an \emph{open} $d$--manifold (never a product with $X$),
by Remark~\ref{rem:noproduct}.  Conversely, when these conditions hold and all
finite subgroups have periodic (e.g.\ cyclic) cohomology, the locally symmetric
orbifold $\Gamma\backslash X$ admits an equivariant rational-homology resolution
(Theorem~\ref{thm:resolution}) producing the desired closed manifold.  Explicit
witnesses are furnished by suitable uniform lattices in $\mathrm{SO}_0(d,1)$
($d\ge5$ odd) all of whose torsion is order-$2$ rotations about codimension-$2$
totally geodesic subspaces (Proposition~\ref{prop:witness}).  The decisive point
distinguishing this from the impossibility of free $2$--group actions on acyclic
spaces is that rational acyclicity carries \emph{no} mod-$2$ Smith obstruction
(Remark~\ref{rem:nosmith}).

\begin{thebibliography}{9}
\bibitem{Lueck}
W.~L\"uck, \emph{Survey on classifying spaces for families of subgroups},
in: Infinite Groups: Geometric, Combinatorial and Dynamical Aspects,
Progr.\ Math.\ \textbf{248}, Birkh\"auser, Basel, 2005, pp.\ 269--322.
\bibitem{Weinberger}
S.~Weinberger, \emph{The Topological Classification of Stratified Spaces},
Chicago Lectures in Mathematics, University of Chicago Press, Chicago, 1994.
\end{thebibliography}

\end{document}
